% !TEX root = ../Projektdokumentation.tex

\section{Analysephase} 
\label{sec:Analysephase}

\subsection{Ist-Analyse} 
\label{sec:IstAnalyse}
Auf Musikfestivals erfolgt der Verkauf von Speisen und Getränken in der Regel ausschließlich vor Ort an einzelnen Verkaufsständen. 
Der typische Ablauf gestaltet sich wie folgt:

\begin{itemize}
    \item Besucher*innen stellen sich vor Ort an einem Stand an
    \item Bestellung wird direkt beim Personal aufgegeben
    \item Bezahlung erfolgt unmittelbar (bar oder per Karte)
    \item Bestellung wird anschließend zubereitet
    \item Besucher*innen warten an der Ausgabe
    \item Übergabe erfolgt direkt am Stand
\end{itemize}

Dieser Prozess findet sequenziell statt, wodurch sich insbesondere zu Stoßzeiten lange Warteschlangen bilden. Das Standpersonal muss 
gleichzeitig Bestellungen entgegennehmen, kassieren, zubereiten und ausgeben, was zu Stresssituationen und erhöhter Fehleranfälligkeit führt.
Besucher*innen haben momentan keine direkte Information über die aktuellen Wartezeiten, die Verfügbarkeit von Produkten oder den Status 
ihrer Bestellung.

\subsection{Soll-Konzept}
\label{sec:SollKonzept}
Zur Behebung der in der Ist-Analyse beschriebenen Schwächen soll ein digitales Bestell- und Abholsystem entwickelt werden, das die bisherigen
manuellen Abläufe auf dem Festival teilweise ersetzt und organisatorisch verbessert.

\subsubsection{Benutzersicht}
Besucher*innen sollen sich über ihre Ticketnummer am System anmelden können. Nach erfolgreicher Anmeldung soll ein Ticket-gebundenes Guthaben
zur Verfügung stehen, das für Bestellungen genutzt werden kann. Außerdem soll dem Benutzer bzw. der Benutzerin ein QR-Code zur Verfügung gestellt werden,
der zur Abholung der Bestellung verwendet wird. Innerhalb der Anwendung sollen Produkte standbezogen ausgewählt und bestellt werden können.
Zusätzlich sollen Informationen über Wartezeiten und Produktverfügbarkeiten angezeigt werden. Ebenso soll der aktuelle Bearbeitungsstatus der Bestellung
jederzeit einsehbar sein.

\subsubsection{Standsicht}
Für das Standpersonal soll eine Verwaltungsoberfläche bereitgestellt werden, in der eingehende Bestellungen übersichtlich dargestellt werden. Die Bestellungen
sollen dort nach Bearbeitungsstatus und gegebenenfalls Priorität verwaltet werden können. Darüber hinaus soll das Personal den Status einzelner Bestellungen ändern
sowie Produktbestände einsehen und bei Bedarf aktualisieren können.

\subsubsection{Bestands- und Bestellverwaltung}
Das System soll Produktbestände zentral verwalten. Bestellungen dürfen nur dann abgeschlossen werden, wenn die benötigten Produkte in ausreichender Menge verfügbar sind.
Nach Abschluss einer Bestellung oder Ausgabe soll der Bestand automatisch angepasst werden.

\subsubsection{Sonderfälle}
Es soll eine Möglichkeit zum stornieren von Bestellungen bereit gestellt werden, darauf folgend muss das zuvor belastete Guthaben automatisch zurückgebucht werden. 
Zusätzlich soll das System eine bevorzugte Behandlung von VIP-Besucher*innen ermöglichen.

\subsubsection{Zielsetzung des Soll-Konzepts}
Durch die Umsetzung des Soll-Konzepts sollen Wartezeiten reduziert, Abläufe an den Verkaufsständen strukturiert und die Transparenz für alle beteiligten Benutzergruppen
erhöht werden. Gleichzeitig soll die digitale Verarbeitung der Bestellungen eine bessere Nachvollziehbarkeit und effizientere Bestandsverwaltung ermöglichen.

\subsection{Wirtschaftlichkeitsanalyse}
\label{sec:Wirtschaftlichkeitsanalyse}
Ziel der Wirtschaftlichkeitsanalyse ist es, die finanzielle Vorteilhaftigkeit der entwickelten Anwendung zu bewerten. 
Da es sich um ein Schulprojekt handelt, basieren die folgenden Berechnungen auf realistischen Annahmen.

\subsubsection{\gqq{Make or Buy}-Entscheidung}
\label{sec:MakeOrBuyEntscheidung}
Im Vorfeld der Entwicklung wurde geprüft, ob eine bestehende Softwarelösung eingesetzt werden kann (Buy) oder eine Eigenentwicklung (Make) sinnvoll ist.
Im Rahmen dieses Projekts wurde die Eigenentwicklung gewählt, da die Anforderungen sehr spezifisch sind und eine flexible Erweiterung des Systems ermöglicht werden soll.

\subsubsection{Projektkosten}
\label{sec:Projektkosten}
Die Projektkosten, die während der Entwicklung des Projektes anfallen, sollen im Folgenden kalkuliert werden. Dafür müssen neben den Personalkosten, die durch die Realisierung 
des Projektes verursacht werden, auch noch die groben Aufwendungen für die Ressourcen berücksichtigt werden. Im Sinne dieses Projektes haben wir nun mit einem Stundensatz von
8€ für einen Schüler und einem Stundensatz von 30€ für einen Lehrer gerechnet.

Die Kosten, die wir für die jeweiligen Vorgänge des Projektes berechnet haben, sowie die gesamten Projektkosten lassen sich der Tabelle~\bhref{tab:Kostenaufstellung} entnehmen.
\tabelle{Kostenaufstellung}{tab:Kostenaufstellung}{Kostenaufstellung} 

\subsubsection{Nutzen / Einsparungen}
\label{sec:Nutzen_Einsparungen}
Annahmen:
\begin{itemize}
    \item 15.000 Besucher pro Tag
    \item 60 \% nutzen das System → 9.000 Bestellungen
    \item 2 € Mehrumsatz pro Bestellung
\end{itemize}

Berechnung:

9.000 × 2 € = 18.000 € Mehrumsatz pro Tag

Zusätzliche wirtschaftliche Vorteile:
\begin{itemize}
    \item Reduzierter Personalbedarf → Einsparung von Personalkosten
    \item Bessere Planbarkeit → weniger Lebensmittelverschwendung
    \item Schnellere Abwicklung → mehr Verkäufe möglich
    \item Höhere Kundenzufriedenheit → langfristige Umsatzsteigerung
\end{itemize}

\subsubsection{Amortisationsdauer}
\label{sec:Amortisationsdauer}
Die Amortisationsdauer beschreibt den Zeitraum, nach dem sich die Investitionskosten durch den generierten Mehrwert ausgeglichen haben.

Bei Gesamtprojektkosten in Höhe von 2.410 € und einem geschätzten täglichen Mehrumsatz von 18.000 € ergibt sich folgende Berechnung:

2.410 € / 18.000 € = 0,134 Tage

Dies entspricht einer Amortisationsdauer von ca. 3,2 Stunden.

Somit amortisiert sich das System bereits innerhalb weniger Stunden am ersten Veranstaltungstag. Selbst bei deutlich geringeren Nutzerzahlen
oder Umsätzen ist von einer sehr schnellen Refinanzierung auszugehen.

\subsection{Anwendungsfälle}
\label{sec:Anwendungsfaelle}
Im Laufe der Analysephase wurde ein Use-Case-Diagramm erstellt, um eine ungefähre Übersicht über die Anwendungsfälle zu erhalten.
Das Diagramm ist im Anhang unter \bhref{app:Use-Case-Diagramm} zu finden und bildet alle Funktionen ab, die aus Endanwendersicht benötigt werden.

\subsection{Qualitätsanforderungen}
\label{sec:Qualitaetsanforderungen}
Die Qualitätsanforderungen konkretisieren die Erwartungen an die Güte des Systems und dienen als Maßstab für die Bewertung der Umsetzung.

\subsubsection{Funktionale Qualität}
Das System muss alle definierten Funktionen vollständig und korrekt abbilden. Bestellungen dürfen nur bei ausreichendem Guthaben und verfügbaren 
Produkten durchgeführt werden. Statusänderungen müssen nachvollziehbar und konsistent erfolgen.

\subsubsection{Benutzerfreundlichkeit}
Die Anwendung soll eine einfache und verständliche Bedienung ermöglichen. Alle relevanten Informationen, wie Bestellstatus oder Wartezeiten, sollen
klar und übersichtlich dargestellt werden.

\subsubsection{Performance}
Die Reaktionszeiten des Systems sollen so gering wie möglich gehalten werden, um eine flüssige Nutzung zu gewährleisten. Insbesondere das Aufgeben
von Bestellungen und das Aktualisieren von Statusinformationen sollen ohne spürbare Verzögerung erfolgen.

\subsubsection{Zuverlässigkeit und Stabilität}
Das System soll auch unter hoher Last stabil funktionieren. Fehlerhafte Zustände, wie doppelte Bestellungen oder falsche Bestände, sollen vermieden werden.

\subsubsection{Nachvollziehbarkeit}
Alle relevanten Aktionen, insbesondere Bestellungen, Stornierungen und Statusänderungen, sollen im System nachvollziehbar gespeichert werden.

\subsubsection{Erweiterbarkeit}
Das System soll so konzipiert sein, dass zukünftige Erweiterungen, wie zusätzliche Benutzerrollen oder neue Funktionen, ohne grundlegende Änderungen möglich sind.

\subsection{Lastenheft/Fachkonzept}
\label{sec:Lastenheft}
Als Ergebnis der Analysephase wurde ein Lastenheft erstellt. Dieses umfasst alle Anforderungen der Auftraggeber an die zu erstellende Anwendung.
Ein Auszug aus dem Lastenheft mit Fokus auf die konkreten Anforderungen befindet sich im Anhang \bhref{app:Lastenheft}.